\BourbakiStarChapter{引言}

代数学本质上关注的是计算，即在集合的元素上施行“代数运算”，其中最著名的例子由初等算术的“四则运算”给出。

这里不是追溯逐步抽象这一缓慢过程的地方；通过这一过程，代数运算的概念最初局限于自然数和可测的量，逐渐扩大其领域，同时“数”的概念得到了平行的推广，直到超越后者，它被应用于不具有“数值”特征的元素，例如集合的置换（见第一章的历史注记）。毫无疑问，正是这些连续扩展的可能性——其中计算的形式保持不变，而经受这些计算的数学实体的性质却变化很大——导致了现代数学指导原则的逐渐确立，即数学实体本身并不重要；重要的是它们之间的关系（见第一卷）。无论如何，可以肯定的是，代数学远在数学的其他分支之前就达到了这种抽象水平，并且长期以来习惯于被视为对代数运算的研究，而独立于这些运算可以应用于的数学实体。

失去了任何特定特征，通常代数运算所依据的共同概念非常简单：对两个元素 \(a, b\) （属于同一集合 \(\mathbf{E}\) ）施行代数运算，意味着将有序对 \((a, b)\) 与一个明确定义的第三元素 \(c\) （属于集合 \(\mathbf{E}\) ）相关联。换言之，这个概念中除了函数之外别无他物：给定一个代数运算，就是给定一个定义在 \(\mathbf{E} \times \mathbf{E}\) 上且取值于 \(\mathbf{E}\) 的函数；其唯一特殊之处在于，该函数的定义域是两个与函数取值所在集合相同的集合的乘积；对于这样的函数，我们称之为合成律。

除了这些“内律”之外，数学家们（主要是在几何学的影响下）还被引导去考虑另一种类型的“合成律”；这就是“作用律”，其中除了集合 E（它仍然可以说是处于前景位置）之外，还有一个辅助集合 \(\Omega\) 出现，其元素称为算子：该律这一次将有序对 \((\alpha, a)\) ——由一个算子 \(\alpha \in \Omega\) 和一个元素 \(a \in E\) 组成——与 E 的第二个元素 b 相关联。例如，欧几里得空间 E 上给定中心的一个位似，将实数 k（“位似比”，此处为算子）和 E 的一点 A 与 E 的另一点 A' 相关联：这就是 E 上的一个作用律。

按照一般定义（《集合论》第四章第1节第4号），在集合 E 上给定一个或多个合成律或作用律，就在 E 上定义了一个结构；对于这样定义的结构，我们专门保留代数结构这一名称，而对它们的研究就构成了代数学。

代数结构有若干种（《集合论》第四章第1节第4号），它们一方面由定义它们的合成律或作用律来刻画，另一方面由这些律所满足的公理来刻画。当然，这些公理并非任意选取的，而恰恰是应用中出现的绝大多数律所具有的性质，例如结合性、交换性等等。第一章主要致力于阐述这些公理以及由它们推出的一般性结论；此外，还对两种最重要的代数结构类型作了更细致的研究，即群（其中只出现一个合成律）和环（有两个合成律），域结构是环的一种特殊情形。

第一章中还给出了带算子群和带算子环的定义，其中除了合成律之外，还出现一个或多个作用律。最重要的带算子群是模，向量空间是模的一个特例，它们在经典几何和现代分析中都起着重要作用。对模结构的研究源于线性方程组的研究，因此得名线性代数；这方面的结果见第二章。

同样地，最常出现的带算子环称为代数（或超复系）。在第三、四章中，对两类特殊的代数作了详细研究：外代数，它与由之导出的行列式理论一起成为线性代数的一个有力工具；以及多项式代数，它是代数方程理论的基础。

第五章阐述交换域的一般理论及其分类。这一理论的起源是含一个未知数的代数方程的研究；当年产生这一理论的问题如今已几乎只有历史意义，但交换域的理论对代数学仍然是基础性的，它一方面是一代数论的基础，另一方面也是代数几何的基础。

由于自然数集有两个合成律，即加法和乘法，经典算术（或数论）——即对自然数的研究——从属于代数学。然而，与由这两个律所定义的代数结构相关联的，还有由序关系“a 整除 b”所定义的结构；而经典算术的独特之处恰恰在于研究这两个相关联结构之间的关系。这并不是序结构与代数结构通过“整除性”关系相关联的唯一例子：这一关系在多项式环中同样起着重要作用。因此，第六章将对此作一般性研究；这一研究将在第七章中应用于确定某些特别简单的环上的模的结构，特别是“初等因子”理论。

第八、九章致力于更专门的理论，但它们在分析中有许多应用：一方面是半单模与半单环的理论，它与群的线性表示理论密切相关；另一方面是二次型与埃尔米特型的理论，以及与之相伴的群的研究。最后，初等几何（仿射、射影、欧几里得等）在第二、六、九章中从纯代数的观点加以研究：这里不过是用一种新的语言来表达别处已经得到的代数学结果，但这种语言特别适合代数几何和微分几何的后续发展，本章正是作为它们的引论。
